\section{Thesis Structure} \label{sec:15-introduction-structure}

\begin{foldable}
  The remainder of this thesis is organised as follows.
  \textbf{Chapter~\ref{ch:20-background}} establishes the background necessary to situate and appreciate the three research contributions.
  It proceeds along two parallel tracks: Track~A covers the foundations of deep learning, model compression, and \ac{KD} theory, then surveys the restricted-access \ac{KD} landscape — black-box, few-shot, and data-free settings — together with the generative and contrastive tools that underpin Chapters~\ref{ch:30-research01} and~\ref{ch:40-research02}; Track~B introduces dataset distillation, influence functions, and \ac{OT}, forming the conceptual basis for Chapter~\ref{ch:50-research03}.
  Readers already familiar with \ac{KD} fundamentals may treat the two tracks as independent references and proceed directly to the contribution chapters.
\end{foldable}

\begin{foldable}
  \textbf{Chapter~\ref{ch:30-research01}} presents DivBFKD, a diversity-aware training scheme for black-box few-shot \ac{KD} on image classification tasks.
  The chapter addresses the problem of limited data availability combined with restricted teacher access, proposing a WGAN-based generation pipeline augmented with an adaptive high-confidence image scheme that broadens the diversity of synthetic training examples without requiring access to teacher internals.
  \textbf{Chapter~\ref{ch:40-research02}} presents DIPKD, which pushes the access constraints further to the extreme black-box data-free setting, where no real training data exists and only top-1 hard labels are available from the teacher.
  DIPKD introduces a three-phase pipeline that synthesises diverse image priors, optimises their diversity through contrastive learning, and recovers soft distillation signals via a primer student — achieving state-of-the-art performance across twelve benchmarks.
  Together, Chapters~\ref{ch:30-research01} and~\ref{ch:40-research02} can be read as a progression along the data-availability axis: from few-shot to fully data-free, while maintaining the black-box teacher assumption throughout.
  \textbf{Chapter~\ref{ch:50-research03}} presents TAKE, which shifts focus from image \ac{KD} to text dataset distillation.
  TAKE addresses a structural bias in influence-function-based sample scoring — the hard-sample bias that causes noisy examples to dominate selection at convergence — by integrating knowledge estimates along the full training trajectory and then selecting a compact, representative corpus via discrete \ac{OT}.
  Although the setting differs from Chapters~\ref{ch:30-research01} and~\ref{ch:40-research02}, the unifying concern is the same: eliciting diverse, high-quality knowledge under practical constraints on data and access.
\end{foldable}

\begin{foldable}
  \textbf{Chapter~\ref{ch:70-conclusion}} synthesises the contributions across all three chapters.
  It draws out the recurring themes — distributional fidelity as a principled design objective, practical restrictions on teacher access and training data, and principled quality control through generative modelling and \ac{OT} — and situates them within the broader trajectory of \ac{KD} research.
  The chapter closes with a discussion of the limitations of each contribution and identifies the most promising directions for future work.
\end{foldable}

\begin{foldable}
  The thesis begins with the background chapter, which may be read in full or selectively by track depending on the reader's familiarity with the relevant literature.
\end{foldable}
